% Solutions to the exercises of Chapter 26.
\begin{solutions}
  \item Take random $H_x \in \R^{8\times39}$, $H_f \in \R^{8\times3}$, an SPD prior $P$ and $\sigma^2 = 0.3$. Update once with $(A\T r, A\T H_x, \sigma^2I_5)$, and once the augmented state
        $(\delta x, \delta p_f)$ with prior $\blkdiag(P, \lambda I_3)$ and $(r, [H_x\ H_f], \sigma^2I_8)$. The $\delta x$ blocks of the two posteriors differ by
        \begin{center}\small
        \begin{tabular}{@{}l c c c c@{}}
        \toprule
        $\lambda$ & $10^2$ & $10^4$ & $10^6$ & $10^8$ \\
        \midrule
        $\max|\Delta P|$ & $5\times10^{-2}$ & $6\times10^{-4}$ & $6\times10^{-6}$ & $6\times10^{-8}$ \\
        $\max|\Delta\delta\hat x|$ & $10^{-2}$ & $10^{-4}$ & $10^{-6}$ & $10^{-8}$ \\
        \bottomrule
        \end{tabular}
        \end{center}
        The difference falls like $1/\lambda$, so in the limit of no prior the two updates coincide, as \cref{thm:nullspace-optimal} states. Beyond $\lambda \approx 10^{10}$ the
        augmented computation itself loses accuracy, through cancellation in $\lambda - \lambda(\cdots)$. This is one more reason to eliminate the feature algebraically.

  \item With $Q = [Q_1\ \ A]$ from the QR decomposition of $H_f$, the first three rotated equations are
        \[
          Q_1\T r = Q_1\T H_x\,\delta x + R_f\,\delta p_f + Q_1\T n, \qquad R_f = Q_1\T H_f \ \text{(upper triangular, invertible)}.
        \]
        For any $\delta x$, they can be satisfied exactly by one $\delta p_f$, which is why they carry no information about $\delta x$. Once the update has produced $\delta\hat x$, they
        determine the feature correction:
        \[
          \delta\hat p_f = R_f\inv Q_1\T\big(r - H_x\,\delta\hat x\big), \qquad \hat p_f \gets \hat p_f + \delta\hat p_f .
        \]
        The noise $Q_1\T n$ is independent of $A\T n$, which drove the update, so the covariance follows directly:
        $P_{ff} = R_f\inv\big(Q_1\T H_xP^+H_x\T Q_1 + \sigma^2I\big)R_f\mT$ and $P_{fx} = -R_f\inv Q_1\T H_xP^+$. Numerically, these agree with the augmented update of Exercise~1 to within $10^{-7}$ (Joseph
        form). This is the back-substitution step of a Schur-complement solver (\cref{sec:schur-ba}). It is what a hybrid filter would do to promote a long-lived
        feature into the state as a landmark, with exactly these covariance blocks.

  \item For two views, $H_f$ is $4\times3$ of rank 3 and its left nullspace is one-dimensional, so $A_j$ is a single unit vector. It can be found without the Jacobian. Let
        $\epsilon(z_1, z_2) = \bar x_2\T E\bar x_1$ be the epipolar function of the two clones' estimated relative pose. For any feature position, its two predicted images satisfy the epipolar
        constraint, so $\epsilon\big(\pi_1(p_f), \pi_2(p_f)\big) \equiv 0$. Differentiating with respect to $p_f$ gives $g\T H_f = 0$ with
        \[
          g = \frac{\partial\epsilon}{\partial(z_1, z_2)} = \big((E\T\bar x_2)_1,\ (E\T\bar x_2)_2,\ (E\bar x_1)_1,\ (E\bar x_1)_2\big), \qquad A_j = g/\|g\| .
        \]
        The projected residual $r_o = g\T r/\|g\|$ is the first-order change of the epipolar residual divided by the norm of its gradient, which is the signed Sampson distance of
        \cref{ch:multiview}. It measures how far the two observations are from satisfying the coplanarity of the two rays and the baseline, and nothing else.
        Errors along the epipolar lines are absorbed by the feature position. For the worked example ($R = I$, $t = (-b, 0, 0)$, $\bar x_1 = (0, 0, 1)$), $g = (0, -b, 0, b)$, which gives
        $A_j = (0, 1, 0, -1)/\sqrt2$ up to sign, as found there. A numerical check at a random configuration, with $H_f$ obtained by finite differences, gives
        $|g\T H_f|/(\|g\|\|H_f\|) = 6\times10^{-11}$.

  \item Applying the three reflections of \cref{alg:nullspace-householder} to $[H_f\ |\ H_x\ |\ r]$ with $m = 2M$ rows and $k = 6M$ clone columns costs about $4(k + 1)(3m - 3)$ flops. Forming
        $A$ from a full QR and multiplying costs about $12m^2$ to build $Q$ plus $2(m - 3)m(k + 1)$ for $A\T[H_x\ r]$:
        \begin{center}\small
        \begin{tabular}{@{}l c c c@{}}
        \toprule
        & $M = 3$ & $M = 6$ & $M = 10$ \\
        \midrule
        Householder in place & 1.1~kflop & 4.9~kflop & 14~kflop \\
        explicit $A$ & 1.1~kflop & 9.7~kflop & 46~kflop \\
        \bottomrule
        \end{tabular}
        \end{center}
        The reflections cost $\mathcal{O}(M^2)$ against $\mathcal{O}(M^3)$, and they need no $m\times m$ workspace. The results agree to $10^{-15}$ up to the sign of each row. A different
        orthonormal basis of the nullspace (from an SVD, say) gives a different $H_o$, but the same $H_o\T H_o$, the same $H_o\T r_o$ and the same gating statistic (checked to all
        digits), because two bases differ by an orthogonal $U$: $A' = AU$.

  \item For features $j \ne k$, $\Cov(n_{o,j}, n_{o,k}) = A_j\T\Cov(n_j, n_k)A_k = 0$ when the raw noises are independent, and $\Cov(n_{o,j}) = \sigma^2A_j\T A_j = \sigma^2I$. The projected
        noises are jointly Gaussian (linear functions of Gaussian noise) and uncorrelated, hence independent, so the stacked noise covariance is $\sigma^2I_d$. The
        projected \emph{residuals} are correlated nevertheless:
        \[
          \Cov(r_{o,j}, r_{o,k}) = H_{o,j}PH_{o,k}\T \ne 0,
        \]
        because every residual depends on the same uncertain state, here on clones that the two features share. These are the off-diagonal blocks of $S = H_oPH_o\T + \sigma^2I$,
        and the stacked update accounts for them exactly. The measurement-noise matrix must not contain them: they are already in $P$, and adding them to $R$
        would count the state uncertainty twice.
\end{solutions}
